\section{Números 0\textendash{}100}
\textbf{Patrón base.} Del 0 al 20 se memorizan; a partir del 21 se construyen ``unidad + en + decena''.

\begin{tabular}{@{}lll@{}}
\toprule
0--10 & \textit{nul, een, twee, drie, vier, vijf, zes, zeven, acht, negen, tien} &  \\
11--20 & \textit{elf, twaalf, dertien, veertien, vijftien, zestien, zeventien, achttien, negentien, twintig} &  \\
Tens & \textit{dertig, veertig, vijftig, zestig, zeventig, tachtig, negentig, honderd} &  \\
\bottomrule
\end{tabular}

\textbf{Construcción 21\textendash{}99}
\begin{itemize}[leftmargin=*]
  \item Forma: unidad + \textit{en} + decena. Ej.: 24 = \textit{vierentwintig}, 57 = \textit{zevenenvijftig}.
  \item Excepción menor: 1 = \textit{een} mantiene la raíz \textit{een} sin acento.
\end{itemize}

\textbf{Mini práctica}
\begin{itemize}[leftmargin=*]
  \item Escribe y lee en voz alta: 13, 28, 46, 73, 99.
  \item Dictado inverso: tu compañero dice el número en neerlandés, escribes la cifra.
\end{itemize}
